\subsection{Norms and traces in etale algebras}

Let A be an etale algebra of (finite) degree n over K. By definition there exist then an extension L of K and distinct homomorphisms \(u_{1}, \ldots, u_{n}\) of A into L with the following properties.

\begin{enumerate}
\def\labelenumi{\alph{enumi})}
\item
  every homomorphism of \(\mathbf{A}\) into \(\mathbf{L}\) is equal to one of the \(u_{i}\) (V, p.~29, Cor.)
\item
  there exists an L-algebra isomorphism \(u: \mathbf{A}_{(L)} \to L^n\) such that
\end{enumerate}

\[
u (1 \otimes x) = \left(u _ {1} (x), \dots , u _ {n} (x)\right) \quad \text { for   all } \quad x \in A.
\]

Moreover, every algebraically closed extension L of K has the above properties (V, p.~30, Prop. 2).

We fix \(L, u_1, \ldots, u_n\) in what follows. Let \(x \in A\) ; we shall prove the formulae

\[
\operatorname{Tr} _ {\mathrm{A} / \mathrm{K}} (x) \cdot 1 = \sum_ {i = 1} ^ {n} u _ {i} (x), \quad \mathrm{N} _ {\mathrm{A} / \mathrm{K}} (x) \cdot 1 = \prod_ {i = 1} ^ {n} u _ {i} (x).\tag{3}
\]

Let \(v\) be multiplication by \(1 \otimes x\) in \(\mathbf{A}_{(L)}\) ; with respect to the basis of \(\mathbf{A}_{(L)}\) which is the image under \(u^{-1}\) of the canonical basis of \(L^n\) , the matrix of the linear mapping \(v\) is diagonal, with diagonal elements \(u_1(x), \ldots, u_n(x)\) . We conclude that

\[
\operatorname{Tr} _ {\mathrm{A} _ {(\mathrm{L})} / \mathrm{L}} (1 \otimes x). 1 = \sum_ {i = 1} ^ {n} u _ {i} (x), \text { whence } \operatorname{Tr} _ {\mathrm{A} / \mathrm{K}} (x). 1 = \sum_ {i = 1} ^ {n} u _ {i} (x) \text { by } (1);
\]

the case of norms is treated similarly.

Further let \((x_{1},\ldots ,x_{n})\) be a sequence of elements of A, let U be the matrix

\[
\left(u _ {i} \left(x _ {j}\right)\right) _ {1 \leqslant i, j \leqslant n}
\]

and let \((t_{ij}) = {}^t\mathrm{U} \cdot U\) . By the first formula (3) we have

\[
\operatorname{Tr} _ {\mathrm{A} / \mathrm{K}} \left(x _ {i} x _ {j}\right) \cdot 1 = \sum_ {k = 1} ^ {n} u _ {k} \left(x _ {i} x _ {j}\right) = \sum_ {k = 1} ^ {n} u _ {k} \left(x _ {i}\right) u _ {k} \left(x _ {j}\right) = t _ {i j};
\]

passing to determinants, we obtain

\[
\mathrm{D} _ {\mathrm{A} / \mathrm{K}} (x _ {1}, \dots , x _ {n}) \cdot 1 = [ \det u _ {i} (x _ {j}) ] ^ {2}.\tag{4}
\]

PROPOSITION 1. --- Let A be a commutative algebra of finite degree over K. Then the following conditions are equivalent :

\begin{enumerate}
\def\labelenumi{\alph{enumi})}
\item
  The algebra \(\mathbf{A}\) is etale.
\item
  There exists a basis of \(\mathbf{A}\) whose discriminant is non-zero.
\item
  For each \(x \neq 0\) in A there exists \(y\) in A such that \(\operatorname{Tr}_{\mathbf{A} / \mathbf{K}}(xy) \neq 0\) .
\end{enumerate}

Further, when these conditions are satisfied, the discriminant of any basis of A is non-zero.

We shall show that when A is assumed etale, the discriminant of A with respect to any basis \((x_{1},\ldots,x_{n})\) of A over K is non-zero; this will in particular establish the implication \(a)\Rightarrow b)\) . By (4), with the above notation, it suffices to show that the matrix U is invertible, or equivalently, that the system of linear equations

\[
\sum_ {i = 1} ^ {n} \lambda_ {i} u _ {i} (x _ {j}) = 0 \quad (\text { for } 1 \leqslant j \leqslant n)\tag{5}
\]

has only the solution \(\lambda_{1} = \cdots = \lambda_{n} = 0\) in L. Now the relation (5) implies \(\sum_{i=1}^{n} \lambda_{i} u_{i}(x) = 0\) for all \(x \in \mathbf{A}\) , whence \(\lambda_{i} = 0\) for \(1 \leqslant i \leqslant n\) , by the theorem on the

linear independence of homomorphisms (V, p.~27, Th. 1).

The equivalence of \(b)\) and \(c)\) is a consequence of the following general lemma:

Lemma 1. --- Let V be a vector space of finite dimension over K and B a bilinear form on V × V. Let \((v_{1}, \ldots, v_{n})\) be a basis of V over K and \(\Delta = \det B(v_{i}, v_{j})\) . Then \(\Delta \neq 0\) if and only if, for each \(x \neq 0\) in V there exists y in V such that \(B(x, y) \neq 0\) .

We have \(\Delta \neq 0\) if and only if the system of linear equations

\[
\sum_ {i = 1} ^ {n} \lambda_ {i} \mathrm{B} (v _ {i}, v _ {j}) = 0 \quad (1 \leqslant j \leqslant n)
\]

has only the solution \(\lambda_1 = \cdots = \lambda_n = 0\) in K. If we put \(x = \sum_{i=1}^{n} \lambda_i v_i\) , the above system is equivalent to B \((x, v_j) = 0\) for \(1 \leqslant j \leqslant n\) , or also, since \((v_1, \ldots, v_n)\) is a basis of V over K, to B \((x, y) = 0\) for all \(y \in V\) , whence the lemma.

Let us show that condition c) implies that A is reduced. Let x be a nilpotent element of A; for any \(y \in A\) the element xy is nilpotent, and so the endomorphism \(L_{xy}\) of the vector space A is nilpotent. Now the following lemma implies that \(\operatorname{Tr}(xy) = 0\) for all \(y \in A\) , whence x = 0 under hypothesis c).

Lemma 2. --- Let V be a vector space of finite dimension over K and u a nilpotent endomorphism of V, then \(\operatorname{Tr}(u)=0\) .

For each integer \(n \geqslant 0\) let \(V_n\) be the image of \(u^n\) . Since \(u\) is nilpotent, there exists an integer \(r \geqslant 0\) such that \(V_0 = V\) , \(V_r = 0\) and \(V_i \neq V_{i+1}\) for \(0 \leqslant i \leqslant r-1\) . Let \(d_i\) be the dimension of \(V_{i-1}\) (for \(1 \leqslant i \leqslant r\) ). There exists a basis \((x_1, \ldots, x_d)\) of \(V\) such that the vectors \(x_j\) with \(d - d_i < j \leqslant d\) form a basis of \(V_{i-1}\) (for \(1 \leqslant i \leqslant r\) ). We have \(u(V_{i-1}) \subset V_i\) and so the diagonal elements of the matrix of \(u\) for the basis \((x_1, \ldots, x_d)\) are zero. Thus we have \(\operatorname{Tr}(u) = 0\) and the lemma follows.

Finally let us show that \(b)\) implies \(a)\) . Let \((x_1, \ldots, x_n)\) be a basis of \(A\) over \(K\) such that \(D_{A/K}(x_1, \ldots, x_n) \neq 0\) . Let \(K'\) be an extension of \(K\) , \(A'\) the \(K'\) -algebra derived from \(A\) by extension of scalars and \(x_i' = 1 \otimes x_i\) for \(1 \leqslant i \leqslant n\) . By Formula (2) (V, p.~47) we have \(D_{A'/K'}(x_1', \ldots, x_n') \neq 0\) . Applying the preceding result to \(A'\) we see that \(A'\) is reduced, hence the algebra \(A\) is etale (V, p.~34, Th. 4).

COROLLARY. --- Let E be an extension of finite degree of K. For E to be separable it is necessary and sufficient that there exist a in E such that \(\operatorname{Tr}_{E/K}(a) \neq 0\) .

The condition is necessary by Prop. 1. Conversely, assume that there exists \(a \in E\) such that \(\operatorname{Tr}_{E/K}(a) \neq 0\) . Given \(x \neq 0\) in \(E\) , if we put \(y = ax^{-1}\) , then \(\operatorname{Tr}_{E/K}(xy) \neq 0\) . Now Prop. 1 shows that \(E\) is an etale algebra over \(K\) , hence a separable extension of \(K\) .

